\SourceLeaf{52}
\[
\begin{array}{lrl}
\log(n-1)\quad .\ .\ .\ .&0.0505666&\\
\log\cos v\quad .\ .\ .\ .\ .&9.6627075&\\
\log\operatorname{tang}h.\quad .\ .\ .\ .&0.0113399&\\
\log\operatorname{cotang}(F-\delta)\quad .\ .&0.0646895&\\\cline{1-2}
\log\cos w\quad .\ .\ .\ .\ .&9.7893035,&\text{unde }w=+52^0\ 0'\ 14''25\\
&&\text{atque }u=10^0\ 36'\ 44''54
\end{array}
\]
\[
\begin{array}{lrl}
\log\operatorname{tang}h.\quad .\ .\ .\ .&0.0113399&\\
\log\cos u\quad .\ .\ .\ .\ .&9.9925074&\\\cline{1-2}
\log\operatorname{tang}G\quad .\ .\ .\ .&0.0188325,&\text{unde }G=46^0\ 14'\ 30''77\\
&&\text{atque }G-\delta=37^0\ 51'\ 47''67
\end{array}
\]
\[
\begin{array}{lrl}
\log\cos G\quad .\ .\ .\ .\ .&9.8398647&\\
\log\operatorname{tang}u.\quad .\ .\ .\ .&9.2726964&\\
\text{Comp. }\log\sin(G-\delta).&0.2119881&\\\cline{1-2}
\log\operatorname{tang}\lambda.\quad .\ .\ .\ .&9.3245492,&\text{unde }\lambda=11^0\ 55'\ 18''31
\end{array}
\]
\[
\begin{array}{lrl}
\log\cos\lambda\quad .\ .\ .\ .\ .&9.9905300&\\
\log\operatorname{cotang}(G-\delta)\quad .\ .&0.1093281&\\\cline{1-2}
\log\operatorname{tang}\varphi\quad .\ .\ .\ .&0.0998581,&\text{unde }\varphi=51^0\ 31'\ 47''19
\end{array}
\]

Haecce poli elevatio paucis tantummodo minutis secundis a vera discrepat. Ex $\lambda$ porro invenitur $\gamma=307^0\ 17'\ 24''91$. sive in tempore $20^{\mathrm h}\ 29^{\mathrm m}\ 9^{\mathrm s}66$, unde praecessio horologii prae tempore sidereo vero $=+10^{\mathrm m}\ 58^{\mathrm s}34$, quae cum determinatione eodem die ex aliis observationibus petita intra dimidiam minuti secundi partem congruit.
\Article{7}
Manca certo foret disquisitio nostra, nisi insuper doceamus, quantus error in determinationem quantitatum incognitarum redundare possit ab erroribus levibus inevitabilibus in altitudinibus mensurandis commissis: simul hinc patebit, quales potissimum stellas eligere conveniat, ut determinationes quam exactissimae inde resultent. Ad hunc finem altitudines $h,h'$ simulque quantitates $\varphi$ et $\lambda$ tamquam variabiles spectemus, quo pacto e differentiatione aequationis $1$ prodit haecce:
\expandafter\def\csname HistoricalNoteText1\endcsname{Azimutha inde a meridie dextram versus per occidentem, septentrionem etc. a $0$ usque ad $360^0$ numerari semper supponimus.}
\SourceLeaf{53}
\[
\cos h\,\mathrm dh=(\sin\delta\cos\varphi-\cos\delta\sin\varphi\cos\lambda)\mathrm d\varphi-\cos\delta\cos\varphi\sin\lambda\,\mathrm d\lambda
\]
Iam constat, si azimuthum circuli verticalis, in quo stella prima observata est, statuatur $=A$, haberi\HistoricalNote{1}
\[
\begin{aligned}
\sin\delta\cos\varphi-\cos\delta\sin\varphi\cos\lambda&=-\cos h\cos A\\
\cos\delta\sin\lambda&=\cos h\sin A
\end{aligned}
\]
unde aequatio praecedens transit in hanc
\[
{}[17]\qquad\mathrm dh=-\cos A\,\mathrm d\varphi-\cos\varphi\sin A\,\mathrm d\lambda
\]
Prorsus simili modo, si azimuthum stellae secundae in observatione altera statuitur $=A'$, prodit e differentiatione aequationis $2$
\[
{}[18]\qquad\mathrm dh'=-\cos A'\,\mathrm d\varphi-\cos\varphi\sin A'\,\mathrm d\lambda
\]
E combinatione harum aequationum demanant sequentes
\[
{}[19]\qquad\mathrm d\varphi=-\frac{\sin A'}{\sin(A'-A)}\mathrm dh+\frac{\sin A\;{}'}{\sin(A'-A)}\mathrm dh'
\]
\[
{}[20]\qquad\cos\varphi\,\mathrm d\lambda=\frac{\cos A'}{\sin(A'-A)}\mathrm dh-\frac{\cos A}{\sin(A'-A)}\mathrm dh'
\]

Hinc sponte patet, errores admodum considerabiles in determinatione tum elevationis poli tum temporis produci posse vel a levissimis erroribus altitudinum observatarum, si $\sin(A'-A)$ fuerit quantitas perparva, i. e. si circulus verticalis secundus vel ipsi primo, vel ei qui primo oppositus est, valde vicinus fuerit: probe itaque cavendum est, ne stellae duae tali modo dispositae combinentur. Simul vero manifestum est, quoties $\sin(A'-A)$ valorem considerabilem positivum seu negativum obtineat, errores observationum certe haud multum auctos in determinatione elevationis poli prodire posse. Idem valet de determinatione temporis, cuius quidem praecisio insuper ab altitudine poli ipsa pendet: hoc vero non methodo sed ipsius rei naturae tribuendum est, quum constet, determinationem temporis minori semper praecisione in latitudinibus maioribus effici, quam in locis aequatori magis vicinis. — Optima quidem semper optio esset, si ad%
\SourceLeaf{54}hibeantur stellae tales, ut $A'-A$ vel $A-A'$ fiat proxime $=90^0$, i. e. si duo circuli verticales proxime sub angulis rectis se intersecent. In hoc casu facile demonstratur, valorem maximum ipsius $\mathrm d\varphi$ fieri $=\BareRoot{(\mathrm dh^2+\mathrm dh'^2)}$, valoremque maximum ipsius $\mathrm d\lambda=\frac{\BareRoot{(\mathrm dh^2+\mathrm dh'^2)}}{\cos\varphi}$. Quodsi itaque de singulis altitudinibus intra $10''$ certi esse possemus, error maximus in determinatione latitudinis foret $=14''$, maximusque error in determinatione anguli $\lambda$, $\frac{14''}{\cos\varphi}$; unde nostris regionibus error maximus in determinatione temporis $=1''5$ temporis. Vix opus erit monere, in hac errorum aestimatione ad solum effectum errorum in observationibus commissorum respici, quod utique sufficit, quum ipsae stellarum maiorum positiones nostra aetate tanta praecisione assignari queant, ut errores pro nihilo habere liceat.

\medskip\centerline{\rule{.25\textwidth}{.3pt}}\medskip
In lectionibus, quas proximo semestri hiberno habiturus sum

\noindent\textit{publice} h. IX theoriam motus cometarum parabolici illustrabo.

\noindent\textit{privatim} h. X \textit{Astronomiam} tradam, audientibus simul copiam exercitationum practicarum tum in observationibus tum in calculo astronomico facturus.
\par\medskip\centerline{\rule{.3\textwidth}{.3pt}}
